\begin{tikzpicture}[
    >=Stealth,
    shorten >=1pt,
    % Thiết lập style để nút nhỏ gọn và nét vẽ đẹp
    state/.style={circle, draw, thick, minimum size=0.8cm, font=\small},
    % Nút 0 là nút khởi tạo, tô đậm viền
    initial/.style={line width=1.5pt},
    every node/.style={font=\scriptsize},
    bend angle=25
]
    % Định nghĩa tọa độ các trạng thái cho giống hình mẫu
    \node[state, initial] (q0) at (0,0) {0};
    \node[state] (q1) at (2.5, -0.6) {1};
    \node[state] (q2) at (2.5, 1.2) {2};
    \node[state, accepting] (q3) at (5, 0) {3/0.7};
    
    % Các đường chuyển trạng thái
    \path[->, thick]
        % Cặp q0 <-> q1 với độ cong tự nhiên
        (q0) edge[bend left=20] node[above, sloped] {a:b/0.1} (q1)
        (q1) edge[bend left=20] node[below, sloped] {a:b/0.2} (q0)
        
        % q1 -> q2 (thẳng đứng)
        (q1) edge node[right] {b:b/0.3} (q2)
        
        % q1 -> q3 và q2 -> q3
        (q1) edge node[below, sloped] {b:b/0.4} (q3)
        (q2) edge node[above, sloped] {a:b/0.5} (q3)
        
        % Vòng lặp tại q3
        (q3) edge[loop right, looseness=5] node[right=2pt] {a:a/0.6} (q3);
\end{tikzpicture}